\chapter{D\'eroulement du stage}
\label{chap:deroulement}

Durant un mois, du \datedebut{} au \datefin{}, notre stage au sein de la
\entreprise{} nous a permis de d\'ecouvrir le fonctionnement d'une structure
publique en charge du pilotage de la recherche scientifique au B\'enin, tout en
menant en parall\`ele notre projet de m\'emoire de fin de formation. Ce chapitre
pr\'esente en premier lieu le minist\`ere de tutelle et la DGRSI, puis d\'ecrit les
activit\'es men\'ees durant le stage. Il se termine par une analyse des apports de
cette exp\'erience et des difficult\'es rencontr\'ees.

\section{Pr\'esentation de la structure}

\subsection{Historique}

La Direction G\'en\'erale de la Recherche Scientifique et de l'Innovation
(DGRSI) est l'une des deux directions g\'en\'erales du Minist\`ere de
l'Enseignement Sup\'erieur et de la Recherche Scientifique (MESRS), aux c\^ot\'es
de la Direction G\'en\'erale de l'Enseignement Sup\'erieur (DGES). Son
organisation et son fonctionnement actuels sont r\'egis par le d\'ecret
n\textsuperscript{o}2023-150 du 12 avril 2023 portant attributions,
organisation et fonctionnement du MESRS, qui pr\'ecise les missions confi\'ees \`a
chacune de ses structures.

\subsection{Raison d'\^etre}

Le MESRS a pour mission de concevoir, mettre en \oe uvre, suivre et \'evaluer la
politique nationale en mati\`ere d'enseignement sup\'erieur, de recherche
scientifique et d'innovation, conform\'ement aux lois en vigueur en R\'epublique
du B\'enin. Au sein de ce minist\`ere, la DGRSI est sp\'ecifiquement charg\'ee de la
conception, de la coordination et du suivi de la mise en \oe uvre de la
politique de l'\'Etat en mati\`ere de recherche scientifique et d'innovation.

\subsection{Objectifs}

Parmi les attributions confi\'ees \`a la DGRSI figurent notamment~: l'\'elaboration
de la politique nationale de la recherche et de l'innovation, en collaboration
avec l'Agence B\'eninoise pour la Recherche et l'Innovation (ABRI)~; la
coordination et la supervision des activit\'es de recherche men\'ees dans les
institutions universitaires publiques et priv\'ees ainsi que dans les instituts,
centres et laboratoires de recherche~; le suivi de ces activit\'es aux plans
global et sectoriel~; la promotion de la formation en sciences et techniques,
notamment \`a l'endroit des femmes et des jeunes~; et l'\'elaboration d'une
plateforme nationale d'informations en science, technologie et innovation.

\section{Organisation technique et manag\'eriale}

\subsection{Structure organisationnelle}

Outre le secr\'etariat de la direction, la DGRSI s'appuie sur des structures
d'accompagnement (Cellule informatique~; Cellule de la Statistique, de la
Documentation et du Pr\'earchivage~; assistant-r\'egisseur) et deux d\'epartements
techniques~: le D\'epartement de la Qualit\'e et du Suivi de la Recherche
scientifique et de l'Innovation, et le D\'epartement du Partenariat et de la
Promotion de la Recherche Scientifique et de l'Innovation.

Nous avons \'et\'e affect\'es \`a la \textbf{Cellule de la Statistique, de la
Documentation et du Pr\'earchivage}, charg\'ee notamment d'\'elaborer les fiches
statistiques sur les activit\'es de la direction, en vue de la collecte et du
traitement des informations, de tenir \`a jour et de pr\'earchiver les documents
de la direction, de produire les statistiques du secteur et de veiller \`a leur
archivage.

\TODOFigure{organigramme_MESRS_DGRSI.png}{Organigramme simplifi\'e du MESRS et de la DGRSI -- la case fonc\'ee indique la structure d'accueil du stage (source~: d\'ecret n\textsuperscript{o}2023-150 du 12 avril 2023)}{fig:organigramme}

\section{Travaux effectu\'es}

Au sein de la Cellule de la Statistique, nous avons particip\'e aux activit\'es
courantes de collecte, de mise \`a jour et de pr\'earchivage des documents et
donn\'ees statistiques de la direction, ce qui nous a permis de d\'ecouvrir
concr\`etement les m\'ecanismes de suivi des activit\'es de recherche et
d'innovation au niveau national.

L'essentiel de notre temps de stage a toutefois \'et\'e consacr\'e \`a la
conception et au d\'eveloppement de notre projet de programmation,
\textbf{CyberGuard B\'enin}, r\'ealis\'e dans le cadre de notre m\'emoire de fin de
formation. Ce travail s'est d\'eroul\'e en plusieurs \'etapes~:

\begin{itemize}[nosep]
  \item clarification du besoin et des objectifs du projet~;
  \item choix de la stack technique (Laravel pour l'API backend, React pour le
    frontend, MySQL comme syst\`eme de gestion de base de donn\'ees)~;
  \item mod\'elisation UML du syst\`eme~: identification des acteurs et des cas
    d'utilisation, diagrammes de classes, d'objets et de s\'equence~;
  \item conception de la base de donn\'ees et d\'efinition du dictionnaire des
    donn\'ees~;
  \item d\'efinition puis red\'efinition compl\`ete de l'identit\'e visuelle de la
    plateforme, suivie de la r\'ealisation de maquettes haute-fid\'elit\'e~;
  \item impl\'ementation du moteur d'analyse de s\'ecurit\'e (sept contr\^oles non
    intrusifs~: HTTPS, certificat SSL, en-t\^etes de s\'ecurit\'e, cookies,
    vuln\'erabilit\'es, fichiers sensibles expos\'es, s\'ecurit\'e email SPF/DMARC),
    du calcul du score et des recommandations~;
  \item impl\'ementation de l'authentification (par email/mot de passe et via
    Google), du tableau de bord, de l'historique des analyses et des
    notifications in-app~;
  \item impl\'ementation du back-office d'administration (gestion des comptes,
    mod\'eration des analyses, param\'etrage du bar\`eme de scoring)~;
  \item tests fonctionnels et corrections successives, jusqu'\`a l'obtention
    d'une plateforme stable et utilisable.
\end{itemize}

\section{Apport du stage sur le plan professionnel}

Ce stage nous a permis de mettre en pratique, sur un projet r\'eel et de bout en
bout, l'ensemble du cycle de vie d'un projet informatique~: de la clarification
du besoin jusqu'\`a une plateforme fonctionnelle, en passant par la
mod\'elisation UML, la conception de base de donn\'ees et le d\'eveloppement
technique. Il a \'egalement renforc\'e nos comp\'etences en d\'eveloppement web
full-stack (Laravel, React), en s\'ecurit\'e applicative (authentification,
protection contre le sondage d'adresses r\'eseau internes, validation des
entr\'ees), ainsi qu'en gestion it\'erative d'un projet~: priorisation des
fonctionnalit\'es, correction de r\'egressions, et adaptation continue de la
plateforme aux retours obtenus au fil du d\'eveloppement. Le passage par la
Cellule de la Statistique nous a par ailleurs sensibilis\'es \`a l'importance
d'une donn\'ee fiable et \`a jour pour le pilotage de politiques publiques -- un
constat qui fait \'echo \`a la probl\'ematique m\^eme de CyberGuard B\'enin~: donner
\`a des acteurs non-techniciens une information fiable et compr\'ehensible sur
laquelle agir.

\section{Difficult\'es rencontr\'ees}

Le d\'eveloppement de CyberGuard B\'enin n'a pas \'et\'e exempt de difficult\'es
techniques. Nous en retenons notamment quatre~:

\begin{itemize}[nosep]
  \item \textbf{Authentification via Google et gestion de session}~:
    l'int\'egration de la connexion via Google (Laravel Socialite) s'est
    heurt\'ee \`a un comportement du middleware de session de Laravel Sanctum,
    qui ne d\'emarre une session que lorsque la requ\^ete est reconnue comme
    provenant du frontend (via l'en-t\^ete \emph{Referer}) -- ce qui n'est pas le
    cas d'un retour depuis Google. Il a fallu forcer explicitement le d\'emarrage
    de session sur les routes concern\'ees.
  \item \textbf{S\'ecurit\'e du moteur d'analyse}~: le moteur devant accepter une
    URL saisie librement par l'utilisateur, une attention particuli\`ere a \'et\'e
    port\'ee \`a emp\^echer qu'il ne soit d\'etourn\'e pour sonder le r\'eseau interne
    du serveur (SSRF)~; la r\'esolution DNS de l'h\^ote est donc v\'erifi\'ee \`a deux
    reprises (\`a la soumission puis juste avant l'analyse effective).
  \item \textbf{Fiabilit\'e du traitement en arri\`ere-plan}~: les analyses \'etant
    trait\'ees de mani\`ere asynchrone via une file d'attente, un oubli de
    d\'emarrage du worker de traitement laissait certaines analyses bloqu\'ees
    ind\'efiniment au statut \og~en cours~\fg{} -- un point de vigilance
    r\'ecurrent en environnement de d\'eveloppement local.
  \item \textbf{Accessibilit\'e des r\'esultats pour un public non-technicien}~:
    conform\'ement \`a l'objectif du projet, les premiers constats techniques
    g\'en\'er\'es (en-t\^etes HTTP manquants, versions de protocole obsol\`etes...)
    se sont r\'ev\'el\'es trop jargonneux pour l'utilisateur cible. L'ensemble des
    descriptions et recommandations a d\^u \^etre reformul\'e en langage courant,
    compr\'ehensible sans connaissance pr\'ealable en s\'ecurit\'e informatique.
\end{itemize}

Ces difficult\'es, une fois identifi\'ees, ont chacune \'et\'e corrig\'ees et
v\'erifi\'ees avant la poursuite du d\'eveloppement.
